\section*{Bab \ref{ch:b3:residues} --- Deret Laurent dan
Teorema Residu}

\begin{solution}{exo:b3:residues:1}
$\frac{\sin z}z = 1 - \frac{z^2}6 + \cdots$: jadi terhapuskan, dengan
\omterm{def:b3:residues:singularities}{residu} $0$. Lalu $\frac{\eu^z - 1}{z^2} = \frac1z + \frac12 +
\frac z6 + \cdots$: kutub sederhana, dengan \omterm{def:b3:residues:singularities}{residu} $1$.
Lalu $\frac1{z(z-1)^2}$: kutub sederhana di $0$ berresidu
$\frac1{(0-1)^2} = 1$; dan kutub rangkap di $1$ berresidu
$\frac{\dd}{\dd z}\bigl(\frac1z\bigr)\big|_{z=1} = -1$.
Lalu $\frac{\cos z}{z^3} = \frac1{z^3} - \frac1{2z} + \cdots$: kutub
berorde $3$, dengan \omterm{def:b3:residues:singularities}{residu} $-\frac12$. Lalu $\eu^{1/z} = \sum_{n\geq0}
\frac{z^{-n}}{n!}$: hakiki, dengan \omterm{def:b3:residues:singularities}{residu} $1$. Dan $\frac1{\sin z}$:
kutub sederhana di $n\pi$; dengan \omterm{def:b3:residues:singularities}{residu} $\frac1{\cos 0} = 1$ di $0$, dan
$\frac1{\cos\pi} = -1$ di $\pi$
(\cref{met:b3:residues:compute}, lewat $g/h'$).
\end{solution}

\begin{solution}{exo:b3:residues:2}
Dengan $z = \eu^{\iu t}$, $\cos t = \frac{z + z^{-1}}2$, dan $\dd t
= \frac{\dd z}{\iu z}$:
\[
\int_0^{2\pi}\frac{\dd t}{a + \cos t}
= \oint_{\abs z = 1}\frac{2\,\dd z}{\iu\,(z^2 + 2az + 1)} .
\]
Akarnya $z_\pm = -a \pm \sqrt{a^2 - 1}$ memenuhi $z_+z_- =
1$ dengan $\abs{z_+} < 1 < \abs{z_-}$; dan \omterm{def:b3:residues:singularities}{residu} di $z_+$ adalah
$\frac{1}{z_+ - z_-} = \frac1{2\sqrt{a^2-1}}$, sehingga integralnya
$\frac2\iu\cdot2\iu\pi\cdot\frac1{2\sqrt{a^2-1}} =
\frac{2\pi}{\sqrt{a^2-1}}$. Saat $a \to 1^+$ ia meledak (sebab
integrannya memuncak di $t = \pi$); sedangkan saat $a \to \infty$ ia berperilaku
seperti $\frac{2\pi}a$, yang cocok dengan $\int\frac{\dd t}a$.
\end{solution}

\begin{solution}{exo:b3:residues:3}
Untuk integral pertamanya: kutub atas $\frac{z^2}{1+z^6}$ di $p =
\eu^{\iu\pi/6}, \iu, \eu^{5\iu\pi/6}$; dan di masing-masingnya,
$\operatorname{Res} = \frac{p^2}{6p^5} = \frac{p^3}{6p^6} =
-\frac{p^3}6$, sedangkan $p^3$ bernilai $\iu, -\iu, \iu$:
sehingga jumlah \omterm{def:b3:residues:singularities}{residunya} $-\frac{\iu}{6}$. Sedangkan busurnya
$O(R^{-4})\cdot O(R) \to 0$:
\[
\int_\R\frac{x^2\,\dd x}{1 + x^6} =
2\iu\pi\Bigl(-\frac \iu6\Bigr) = \frac\pi3 .
\]
Untuk yang kedua: kutub rangkap di $\iu$ milik $\frac1{(1+z^2)^2} =
\frac1{(z-\iu)^2(z+\iu)^2}$:
\[
\operatorname{Res} = \frac{\dd}{\dd z}\,(z +
\iu)^{-2}\Big|_{z=\iu} = \frac{-2}{(2\iu)^3} =
\frac{-2}{-8\iu} = -\frac\iu4,
\qquad
\int_\R\frac{\dd x}{(1+x^2)^2} =
2\iu\pi\cdot\Bigl(-\frac\iu4\Bigr) = \frac\pi2 ,
\]
yang konsisten dengan \cref{exo:b3:fouriertransform:3}(b).
\end{solution}

\begin{solution}{exo:b3:residues:4}
Tutuplah $\frac{\eu^{\iu tz}}{z^2 + a^2}$ di setengah bidang atasnya
(dengan $t \geq 0$): sebab di sana $\abs{\eu^{\iu tz}} =
\eu^{-t\operatorname{Im}z} \leq 1$, sehingga busurnya menyumbang
$O(R^{-2})\cdot O(R) \to 0$. Sedangkan satu-satunya kutub terlingkupinya $\iu a$
bersifat sederhana dengan \omterm{def:b3:residues:singularities}{residu} $\frac{\eu^{-at}}{2\iu a}$:
\[
\int_\R\frac{\eu^{\iu tx}}{x^2 + a^2}\dd x =
\frac{\pi}{a}\,\eu^{-at},
\qquad\text{sehingga}\qquad
\int_\R\frac{\cos(tx)}{x^2+a^2}\dd x = \frac\pi
a\,\eu^{-a\abs t}
\]
(lewat bagian realnya; dan genap terhadap $t$). Inilah pasangan pembalikan
$\widehat{\eu^{-a\abs x}} = \frac{2a}{a^2+\xi^2}$
(\cref{exo:b3:fouriertransform:1}): sehingga kedua perhitungannya
saling membenarkan lewat
\cref{thm:b3:fouriertransform:inversion}.
\end{solution}

\begin{solution}{exo:b3:residues:5}
Untuk pecahan parsialnya: $f = \frac1{z-2} - \frac1{z-1}$.
Pada $\abs z < 1$ (lewat Taylor):
$f = \sum_{n\geq0}\bigl(1 - 2^{-n-1}\bigr)z^n$.
Pada $1 < \abs z < 2$: $\frac1{z-2} =
-\sum_{n\geq0}\frac{z^n}{2^{n+1}}$ dan $-\frac1{z-1} =
-\sum_{n\geq1}z^{-n}$: yakni deret dua sisi yang sejati.
Sedangkan pada $\abs z > 2$: $f = \sum_{n\geq1}\bigl(2^{n-1} -
1\bigr)z^{-n}$. Ketiganya berbeda sebab uraian Laurent
melekat pada \emph{anulus}, bukan pada titik: sehingga setiap daerahnya mempunyai
uraian geometrinya sendiri. Sedangkan koefisien $z^{-1}$-nya (yakni $-1$ dan
$0$ berturut-turut) merupakan integral atas lingkaran yang melingkari
\emph{kesingularan di dalamnya}, bukan \omterm{def:b3:residues:singularities}{residu} di $0$ (sebab $f$
\omterm{def:b3:holomorphic:holo}{holomorfik} di $0$): jadi untuk $1 < \abs z < 2$ koefisiennya
$-1$ adalah $\operatorname{Res}(f, 1)$; sedangkan untuk $\abs z > 2$
koefisiennya $0$ adalah $\operatorname{Res}(f,1) +
\operatorname{Res}(f,2) = -1 + 1$. Adapun \omterm{def:b3:residues:singularities}{residu} $f$: $-1$
di $1$ dan $+1$ di $2$.
\end{solution}

\begin{solution}{exo:b3:residues:6}
(a) \omterm{thm:b3:residues:laurent}{Deret Laurentnya} $\sum_nz^{-n}/n!$ bersuku negatif tak
berhingga banyak: jadi hakiki. Lalu dengan menyelesaikan $\eu^{1/z} = w$ (dengan $w \neq
0$): $\frac1z = \log\abs w + \iu\arg w + 2\iu\pi k$, sehingga
\[
z_k = \frac1{\log\abs w + \iu\arg w + 2\iu\pi k}
\xrightarrow[k\to\infty]{} 0 :
\]
sehingga setiap nilai tak nolnya tercapai tak berhingga sering di dekat $0$ ---
jadi lebih kuat daripada \omterm{ex:b3:lebesgue:gamma}{kepadatan}.
(b) Sepanjang $z = 1/x$ dengan $x \to +\infty$: $\eu^x \to \infty$;
sedangkan sepanjang $z = \iu/y$: modulusnya $1$. Sebuah kutub menuntut
$\abs{f(z)} \to \infty$ \emph{sepanjang setiap pendekatan}: sedangkan di sini
limitnya sekadar tak ada, bahkan di $[0, +\infty]$.
\end{solution}

\begin{solution}{exo:b3:residues:7}
(a) Pada $\abs z = 1$: $\abs{z^7 + z - 1} \leq 3 < 4 =
\abs{-4z^3}$. Lalu Rouché dengan $f = -4z^3$ dan $g = z^7 + z - 1$:
memberikan tiga akar di cakramnya.
(b) Pada $\abs z = 1$: $\abs{z^4 + 1} \leq 2 < 5 = \abs{5z}$:
sehingga satu akar di $\abs z < 1$. Sedangkan pada $\abs z = 2$: $\abs{5z + 1} \leq
11 < 16 = \abs{z^4}$: sehingga empat akar di $\abs z < 2$. Karena itu tiga
akar di anulusnya.
(c) Untuk $P = z^n + a_{n-1}z^{n-1} + \dots$: pada $\abs z = R >
1 + \sum\abs{a_k}$, $\abs{P - z^n} \leq
\bigl(\sum\abs{a_k}\bigr)R^{n-1} < R^n = \abs{z^n}$: sehingga $P$ berakar
tepat $n$ kali di $D(0, R)$ --- yakni d'Alembert--Gauss beserta
kelipatannya, lewat pencacahan belaka.
\end{solution}

\begin{solution}{exo:b3:residues:8}
(a) Andaikan $f(a) = 0$ dengan $f \not\equiv 0$: pilihlah $r$ yang $f$-nya
bebas akar pada lingkaran $C = \partial D(a, r)$ (menurut \omterm{thm:b3:holomorphic:identity}{akar
terpencilnya}) dan $m = \min_C\abs f > 0$. Lalu menurut
\cref{thm:b3:holomorphic:weierstrassconv}, $f_n \to f$ dan
$f_n' \to f'$ secara seragam pada $C$; sehingga untuk $n$ yang besar, $\abs{f_n}
\geq m/2$ pada $C$, jadi
\[
\frac1{2\iu\pi}\int_C\frac{f_n'}{f_n}
\longrightarrow \frac1{2\iu\pi}\int_C\frac{f'}{f} \geq 1
\]
(sebab limitnya mencacah akar $a$; dan kekonvergenannya karena
pembilangnya konvergen seragam sedangkan penyebutnya terbatas
seragam dari bawah). Sedangkan ruas kirinya bilangan bulat yang mencacah akar
$f_n$ di cakramnya: sehingga akhirnya ia haruslah $\geq 1$ ---
yang bertentangan dengan kebebasan akarnya. Jadi $f$ bebas akar.

(b) Tetapkan $w \in \Omega$ lalu terapkan (a) pada \omterm{def:b3:topology:topology}{himpunan terbuka} \omterm{def:b3:topology:connected}{terhubung}
$\Omega\setminus\{w\}$ (sebab membuang sebuah titik dari himpunan bagian terbuka
\omterm{def:b3:topology:connected}{terhubung} $\C$ mengawetkan keterhubungannya) untuk $g_n(z) = f_n(z) -
f_n(w)$, yang bebas akar di sana menurut keinjektifannya, dan konvergen ke $g = f
- f(w)$. Jika $f$ tak konstan, maka $g \not\equiv 0$ pada
$\Omega\setminus\{w\}$, sehingga $g$ bebas akar di sana: jadi $f(z) \neq
f(w)$ untuk setiap $z \neq w$. Dan karena $w$ sembarang, $f$ bersifat
injektif.
\end{solution}

\begin{solution}{exo:b3:residues:9}
Batas juringnya terdiri atas $[0, R]$, busur $A_R$, dan
sinar $\eu^{2\iu\pi/n}[0, R]$ yang dibalik. Di dalamnya terletak
satu-satunya kutub $p = \eu^{\iu\pi/n}$ milik $\frac1{1+z^n}$, dengan
\omterm{def:b3:residues:singularities}{residu} $\frac1{np^{n-1}} = \frac{p}{np^n} = -\frac pn$. Pada
sinar kembalinya, $z = \eu^{2\iu\pi/n}x$ memberikan $z^n = x^n$ dan
$\dd z = \eu^{2\iu\pi/n}\dd x$; sedangkan busurnya
$O(R^{-n})\cdot O(R) \to 0$. Karena itu
\[
\bigl(1 - \eu^{2\iu\pi/n}\bigr)
\int_0^\infty\frac{\dd x}{1 + x^n}
= 2\iu\pi\Bigl(-\frac{\eu^{\iu\pi/n}}n\Bigr),
\quad\text{sehingga}\quad
\int_0^\infty\frac{\dd x}{1+x^n}
= \frac{2\iu\pi}{n\,\bigl(\eu^{\iu\pi/n} -
\eu^{-\iu\pi/n}\bigr)} = \frac{\pi}{n\sin(\pi/n)} .
\]
Untuk $n = 2$: $\frac\pi{2\sin(\pi/2)} = \frac\pi2 =
[\arctan]_0^\infty$. Sedangkan saat $n \to \infty$: nilainya menuju
$1$, dan sungguh integrannya menuju $\mathbf 1_{\intco01}$
(dengan dominasi $\min(1, x^{-2})$ bagi kekonvergenan terdominasinya untuk $n \geq 2$).
\end{solution}

\begin{solution}{exo:b3:residues:10}
Pada $\abs z = R$ yang besar, $\abs f \leq C R^{-2}$: sehingga $\abs{\oint_{C_R}
f} \leq 2\pi R\cdot CR^{-2} \to 0$. Padahal untuk $R$ di luar semua
kutubnya, teorema \omterm{def:b3:residues:singularities}{residunya} memberikan $\oint_{C_R}f =
2\iu\pi\sum_{\text{semua }p}\operatorname{Res}(f, p)$: sehingga jumlah
totalnya lenyap. Untuk $f = \frac1{z(z-1)(z-2)}$: \omterm{def:b3:residues:singularities}{residunya}
$\frac1{(-1)(-2)} = \frac12$ di $0$, $\frac1{1\cdot(-1)} = -1$
di $1$, dan $\frac1{2\cdot1} = \frac12$ di $2$ --- yang berjumlah $0$
seperti diramalkan, dan
\[
\frac1{z(z-1)(z-2)} = \frac{1/2}{z} - \frac1{z - 1} +
\frac{1/2}{z-2} :
\]
sehingga \omterm{def:b3:residues:singularities}{residunya} \emph{adalah} koefisien pecahan parsialnya, dan
identitas berjumlah nol itu menyediakan pemeriksaan kekonsistenan yang cuma-cuma (atau
menentukan koefisien terakhirnya dari yang lain).
\end{solution}

\begin{solution}{exo:b3:residues:11}
Pada lubang kuncinya, dengan penentuan yang dipilih:
tepat di atas irisannya, $\log z = \ln x$; sedangkan tepat di bawahnya, $\log z =
\ln x + 2\iu\pi$. Lalu keempat kepingannya memberikan
\[
\Bigl(1 - \eu^{2\iu\pi(a-1)}\Bigr)\int_\varepsilon^R
\frac{x^{a-1}}{1+x}\dd x + \int_{C_R} + \int_{C_\varepsilon}
= 2\iu\pi\operatorname{Res}(f, -1) .
\]
Untuk busurnya: $\abs{f} \leq \frac{R^{a-1}}{R - 1}$ pada $C_R$ yang
berpanjang $2\pi R$: sehingga sumbangannya $O(R^{a-1}) \to 0$ (sebab $a < 1$);
sedangkan $\abs f \leq \frac{\varepsilon^{a-1}}{1 - \varepsilon}$ pada
$C_\varepsilon$ yang berpanjang $2\pi\varepsilon$: jadi $O(\varepsilon^a)
\to 0$ (sebab $a > 0$). Untuk \omterm{def:b3:residues:singularities}{residunya}: di $z = -1 = \eu^{\iu\pi}$,
$\operatorname{Res} = \eu^{(a-1)\iu\pi}$. Karena itu
\[
I(a) = \frac{2\iu\pi\,\eu^{\iu\pi(a-1)}}{1 -
\eu^{2\iu\pi(a-1)}}
= \frac{2\iu\pi}{\eu^{-\iu\pi(a-1)} - \eu^{\iu\pi(a-1)}}
= \frac{\pi}{-\sin(\pi(a-1))} = \frac{\pi}{\sin\pi a} .
\]
Untuk pencerminan Gammanya: $B(a, 1-a) =
\int_0^1t^{a-1}(1-t)^{-a}\dd t$; lalu substitusi $t =
\frac x{1+x}$ dengan $1 - t = \frac1{1+x}$ dan $\dd t =
\frac{\dd x}{(1+x)^2}$ mengubahnya menjadi $\int_0^\infty
\frac{x^{a-1}}{1+x}\dd x = I(a)$, dan rumus Euler
$B(a, 1-a) = \Gamma(a)\Gamma(1-a)/\Gamma(1)$
(\cref{pb:b3:lebesgue:1}) memberikan
$\Gamma(a)\Gamma(1-a) = \frac\pi{\sin\pi a}$ --- yang
khususnya memberikan $\Gamma(\tfrac12) = \sqrt\pi$ sekali lagi.

\end{solution}

\begin{solution}{exo:b3:residues:12}
(a) Pada $\abs z = 1$: $\abs{z^4} = 1 < 7 \leq \abs{8z + 1}$
(sebab $\abs{8z} - 1 = 7$): sehingga $P$ berakar di $\mathbb D$ sebanyak
$8z + 1$, yakni \emph{satu} (di $-\frac18$).

(b) Pada $\abs z = 3$: $\abs{8z + 1} \leq 25 < 81 =
\abs{z^4}$: sehingga Rouché terhadap $z^4$ memberikan keempat akarnya di
$\abs z < 3$, jadi $4 - 1 = 3$ akar di anulus $1 <
\abs z < 3$. Untuk penajamannya: sebuah akar dengan $\abs z = r \geq 2.1$
akan memenuhi $r^4 = \abs{8z + 1} \leq 8r + 1$, padahal $r^4 -
8r - 1$ naik untuk $r \geq 2$ dan bernilai $19.45 -
16.8 - 1 = 1.65 > 0$ di $r = 2.1$: jadi mustahil. Sehingga ketiga
akar luarnya terletak di $1 < \abs z < 2.1$. (Secara numerik: sebuah akar
real di dekat $-1.95$ dan sepasang sekawan di dekat $1.04 \pm
1.73\iu$ yang bermodulus $2.02$ --- itulah mengapa percobaan
Rouché pada jari-jari tepat $2$ pasti gagal: sebab teoremanya menuntut
dominasi yang sejati, sedangkan akarnya duduk tepat di luar.)

(c) Tak ada akar pada $\iu\R$: sebab $\operatorname{Re}P(\iu t) = t^4 +
1 \geq 1$. Untuk akar di setengah bidang kanannya: pakailah asas
argumen pada batas setengah cakram $\{\abs z \leq
R,\ \operatorname{Re}z \geq 0\}$. Pada busur besarnya, $\arg P
\approx \arg z^4$ berputar sebesar $4\cdot\pi = 2\pi\cdot2$ (sebab
busurnya merentang sudut $\pi$). Sedangkan sepanjang sumbu imajinernya dari
$\iu R$ turun ke $-\iu R$: $P(\iu t) = (t^4 + 1) + 8\iu t$
tetap di setengah bidang kanannya (sebab $\operatorname{Re} > 0$), sehingga
$\arg P$ beragam di dalam $\intoo{-\pi/2}{\pi/2}$ dan kembali
dengan perubahan neto $\to 0$ saat $R \to \infty$ (sebab ujungnya sama-sama
$\approx \arg t^4 = 0$). Jadi lilitan totalnya: $\frac{4\pi + 0}
{2\pi} = 2$: yakni \emph{dua} akar di setengah bidang kanannya ---
yang konsisten dengan numeriknya: sebab sepasang sekawan di $\approx
1.04 \pm 1.73\iu$ berbagian real positif, sedangkan akar realnya
$\approx -0.125$ dan $\approx -1.96$ negatif.
\end{solution}

\begin{solution}{pb:b3:residues:1}
\textbf{1.} Fungsi $\sin(\pi z)$ berakar sederhana tepat di $\Z$
(sebab $\sin\pi z = 0$ jika dan hanya jika $z \in \Z$, dan $(\sin\pi z)' = \pi\cos\pi
z \neq 0$ di sana), dan $\cos(\pi n) \neq 0$: sehingga $\pi\cot(\pi z) =
\pi\cos(\pi z)/\sin(\pi z)$ berkutub sederhana di $\Z$ dengan
\[
\operatorname{Res}(\pi\cot\pi z,\ n)
= \frac{\pi\cos(\pi n)}{\pi\cos(\pi n)} = 1
\]
(lewat aturan $g/h'$, \cref{met:b3:residues:compute}).

\textbf{2.} Fungsi $\frac{\sin(\pi z)}{\pi z} = 1 - \frac{(\pi z)^2}6
+ \frac{(\pi z)^4}{120} - \cdots$ bersifat \omterm{def:b3:holomorphic:holo}{holomorfik} dan tak nol
di dekat $0$: sehingga kebalikannya \omterm{def:b3:holomorphic:holo}{holomorfik}
(\cref{def:b3:holomorphic:holo}: lewat hasil baginya), dengan deret $1 +
\frac{(\pi z)^2}{6} + \frac{7(\pi z)^4}{360} + \cdots$
(kenalilah: yakni koefisien bergaya $(1 - u)^{-1}$ dari $u =
\frac{(\pi z)^2}6 - \frac{(\pi z)^4}{120}$: sehingga koefisien $z^4$-nya
adalah $\frac1{36} - \frac1{120} = \frac{7}{360}$).
Lalu kalikan dengan $\cos(\pi z) = 1 - \frac{(\pi z)^2}2 + \frac{(\pi
z)^4}{24} - \cdots$ lalu bagilah dengan $z$:
\[
\pi\cot(\pi z) = \frac1z\Bigl[1 + \pi^2z^2\Bigl(\frac16 -
\frac12\Bigr) + \pi^4z^4\Bigl(\frac7{360} - \frac1{12} +
\frac1{24}\Bigr)\Bigr] + \cdots
= \frac1z - \frac{\pi^2}3\,z - \frac{\pi^4}{45}\,z^3 -
\cdots
\]
(sebab $\frac7{360} - \frac{30}{360} + \frac{15}{360} =
-\frac8{360} = -\frac1{45}$).

\textbf{3.} Untuk sisi tegaknya $z = \pm(N + \frac12) + \iu y$:
menurut keperiodikan-$\pi$ $\cot$, $\cot(\pi z) = \cot(\pm\frac\pi2
+ \iu\pi y) = -\tan(\iu\pi y) = -\iu\tanh(\pi y)$, yang bermodulus
$\leq 1$. Sedangkan untuk sisi mendatarnya $z = x \pm \iu(N + \frac12)$: dari
$\abs{\cot(a + \iu b)}^2 = \frac{\cos^2a +
\sinh^2b}{\sin^2a + \sinh^2b} \leq \frac{1 +
\sinh^2b}{\sinh^2b} = \coth^2 b$,
\[
\abs{\cot(\pi z)} \leq \coth\bigl(\pi(N + \tfrac12)\bigr)
\leq \coth(\pi/2) < 1.1 .
\]
Jadi kedua batasnya $\leq 2$.

\textbf{4.} Terapkan \cref{thm:b3:residues:residue} pada $F(z) =
\pi\cot(\pi z)f(z)$ di $C_N$ (dengan $N$ di luar semua kutub $f$):
\[
\frac1{2\iu\pi}\oint_{C_N}F = \sum_{n=-N}^{N}f(n) +
\sum_p\operatorname{Res}(F, p),
\]
dengan kutub bulatnya menyumbang $f(n)$ (pertanyaan 1; sebab $f$
\omterm{def:b3:holomorphic:holo}{holomorfik} di sana). Pada $C_N$: $\abs F \leq 2\pi\cdot
C\abs z^{-2} \leq 2\pi C N^{-2}$, dan kelilingnya $8(N +
\frac12)$: sehingga integralnya $O(1/N) \to 0$. Lalu biarkan $N \to
\infty$: diperolehlah rumus penjumlahan yang ditampilkan itu.

\textbf{5.} Peluruhan $\abs f = O(\abs z^{-2})$ membunuh
\omterm{def:b3:holomorphic:contour}{integral konturnya} \emph{sekaligus} membuat $\sum\abs{f(n)}$ konvergen.
Untuk $f(z) = \frac1{z + \frac12}$: jumlah setangkupnya
$\sum_{-N}^N\frac1{n + \frac12}$ berteleskop menjadi $0$ (sebab suku
$n$ dan $-n - 1$ saling meniadakan), dan ruas kanannya adalah
$-\operatorname{Res}\bigl(\frac{\pi\cot\pi z}{z + \frac12},
-\frac12\bigr) = -\pi\cot(-\frac\pi2) = 0$: jadi konsisten ---
padahal $\sum\abs{f(n)}$ divergen; sehingga metodenya menghitung
limit setangkupnya (yakni nilai pokoknya) saja.

\textbf{6.} Untuk $g(z) = \frac{\pi\cot(\pi z)}{z^2}$: kutubnya di
bilangan bulat tak nol dengan \omterm{def:b3:residues:singularities}{residu} $\frac1{n^2}$, dan di $0$
tempat, menurut pertanyaan 2,
\[
g(z) = \frac1{z^3} - \frac{\pi^2}{3z} - \frac{\pi^4}{45}z -
\cdots :
\qquad \operatorname{Res}(g, 0) = -\frac{\pi^2}3 .
\]
Sedangkan \omterm{def:b3:holomorphic:contour}{integral konturnya} atas $C_N$ menuju $0$ seperti pada pertanyaan 4
(sebab $\abs{g} = O(N^{-2})$ pada $C_N$). Karena itu $0 =
\sum_{n\neq0}\frac1{n^2} - \frac{\pi^2}3$: sehingga $2\zeta(2) =
\frac{\pi^2}3$ dan $\zeta(2) = \frac{\pi^2}6$.

\textbf{7.} Untuk $g(z) = \frac{\pi\cot(\pi z)}{z^4} = \frac1{z^5}
- \frac{\pi^2}{3z^3} - \frac{\pi^4}{45z} - \cdots$: \omterm{def:b3:residues:singularities}{residunya}
di $0$ sama dengan $-\frac{\pi^4}{45}$, dan $0 =
2\zeta(4) - \frac{\pi^4}{45}$: sehingga $\zeta(4) =
\frac{\pi^4}{90}$.

\textbf{8.} Dengan $g = \pi\cot(\pi z)/z^{2k}$: \omterm{def:b3:residues:singularities}{residunya} di
$0$ adalah koefisien $a_{2k-1}$ milik $z^{2k-1}$ pada
uraian $\pi\cot(\pi z)$, dan konturnya yang lenyap memberikan
$2\zeta(2k) + a_{2k-1} = 0$: sehingga $\zeta(2k) = -a_{2k-1}/2$, yakni
kelipatan rasional $\pi^{2k}$ sebab koefisien
kotangennya demikian. Lalu satu langkah pembagian lagi melahirkan $a_5 =
-\frac{2\pi^6}{945}$, sehingga $\zeta(6) = \frac{\pi^6}{945}$.
Untuk $\zeta(3)$: kernel alaminya $g = \pi\cot(\pi z)/z^3$
menghasilkan $\sum_{n\neq0}\frac1{n^3} = 0$ menurut kegasalannya --- sehingga
metodenya membuktikan $0 = 0$ dan secara struktural buta terhadap nilai zeta
gasalnya (sebab tak ada bentuk tertutup $\zeta(3)$ yang dikenal; dan
keirasionalannya, menurut Apéry 1978, memerlukan gagasan yang sama sekali
berbeda).

\textbf{9.} Fungsi $F(z) = \frac{\pi\cot(\pi z)}{(z - w)(z + w)}$
berkutub di bilangan bulatnya (dengan \omterm{def:b3:residues:singularities}{residu} $\frac1{n^2 - w^2}$,
sambil mencatat tandanya: $f(n) = \frac1{(n-w)(n+w)} = \frac1{n^2 -
w^2}$) dan berkutub sederhana di $\pm w$ dengan \omterm{def:b3:residues:singularities}{residu}
$\frac{\pi\cot(\pm\pi w)}{\pm2w} = \frac{\pi\cot(\pi
w)}{2w}$ masing-masing (sebab $\cot$ gasal). Lalu argumen pertanyaan 4
(dengan $\abs f = O(\abs z^{-2})$) memberikan
\[
\sum_{n\in\Z}\frac1{n^2 - w^2} + \frac{\pi\cot(\pi w)}{w} =
0,
\qquad\text{yakni}\qquad
\pi\cot(\pi w) = \frac1w + \sum_{n\geq1}\frac{2w}{w^2 -
n^2},
\]
(sebab suku $n = 0$-nya adalah $-\frac1{w^2}$; lalu kelompokkan ulang $\pm n$). Untuk kekonvergenan
normalnya pada \omterm{def:b3:topology:compact}{kompak} $\C\setminus\Z$: bagi $\abs w \leq
R$ dan $n \geq 2R$, $\abs{\frac{2w}{w^2 - n^2}} \leq
\frac{2R}{n^2 - R^2} \leq \frac{8R}{3n^2}$.

\textbf{10.} Untuk $\abs w \leq r < 1$: $\frac{2w}{w^2 - n^2} =
-\frac{2w}{n^2}\cdot\frac1{1 - w^2/n^2} =
-2\sum_{k\geq0}\frac{w^{2k+1}}{n^{2k+2}}$, dengan
$\abs{\text{sukunya}} \leq 2r^{2k+1}/n^{2k+2}$, yang terjumlahkan atas
$(n, k)$: sehingga Fubini untuk deret menyusun ulangnya menjadi
\[
\pi\cot(\pi w) = \frac1w -
2\sum_{k\geq0}\zeta(2k+2)\,w^{2k+1} .
\]
Lalu dengan mencocokkannya terhadap pertanyaan 2: $-2\zeta(2) = -\frac{\pi^2}3$ dan
$-2\zeta(4) = -\frac{\pi^4}{45}$ --- yakni nilai yang sama. Jadi tiga
jalan menuju $\frac{\pi^2}6$: \omterm{thm:b3:hilbert:parseval}{Parseval} (lewat deret Fourier),
trace operator Green dawainya, dan kedua uraian
kotangennya; sedangkan bahwa sebuah jumlah atas frekuensi, sebuah trace operator,
dan sebuah \omterm{def:b3:holomorphic:contour}{integral kontur} bersesuaian bukanlah kebetulan --- sebab masing-masingnya
wajah identitas spektral yang sama.

\textbf{11.} Tetapkan $R \geq 1$ lalu misalkan $n_0$ bilangan bulat terkecil
yang $\geq 2R$. Untuk $\abs z \leq R$ dan $n \geq n_0$:
$\abs{z^2/n^2} \leq \frac14$, sehingga $1 - z^2/n^2 \in
\bar D(1, \frac14)$, tempat logaritma pokoknya bersifat
\omterm{def:b3:holomorphic:holo}{holomorfik}, dan
\[
\Bigl|\log\Bigl(1 - \frac{z^2}{n^2}\Bigr)\Bigr|
\leq \sum_{k\geq1}\frac1k\,\Bigl|\frac{z^2}{n^2}\Bigr|^k
\leq \frac{\abs{z^2/n^2}}{1 - \abs{z^2/n^2}}
\leq \frac{2R^2}{n^2} :
\]
sehingga jumlahnya $S(z) = \sum_{n\geq n_0}\log(1 - z^2/n^2)$ konvergen
normal pada $\bar D(0, R)$, dengan jumlah parsial $S_N$ yang \omterm{def:b3:holomorphic:holo}{holomorfik}
dan $\abs{S_N} \leq 2R^2\zeta(2)$ secara seragam. Lalu karena
$\abs{\eu^a - \eu^b} \leq
\eu^{\max(\abs a,\abs b)}\abs{a - b}$ (lewat batas nilai rata-rata pada
ruasnya), hasil kali ekornya $\prod_{n_0\leq n\leq N} =
\eu^{S_N}$ konvergen seragam pada $\bar D(0, R)$ ke
$\eu^S$ yang bebas akar. Lalu dengan mengalikannya dengan polinomial tetap
$z\prod_{n<n_0}(1 - z^2/n^2)$: $P_N \to P$ secara seragam pada
$\bar D(0, R)$, dan
\cref{thm:b3:holomorphic:weierstrassconv} membuat $P$
\omterm{def:b3:holomorphic:holo}{holomorfik} di sana; dan karena $R$ sembarang, $P$ bersifat utuh. Pada
$\bar D(0, R)$ akar $P$ adalah akar prafaktor
polinomialnya --- yakni bilangan bulat bermodulus $\leq R$, yang masing-masing sederhana
(sebab $(1 - z/n)(1 + z/n)$ berakar sederhana yang berbeda, sedangkan $\eu^S$
tak berakar): jadi himpunan akar $P$ adalah $\Z$, dengan semua akarnya sederhana.

\textbf{12.} Lewat pendiferensialan logaritmik hasil kali
berhingganya, jauh dari akarnya:
\[
\frac{P_N'(z)}{P_N(z)} = \frac1z +
\sum_{n=1}^{N}\frac{-2z/n^2}{1 - z^2/n^2}
= \frac1z + \sum_{n=1}^{N}\frac{2z}{z^2 - n^2} .
\]
Pada sebuah \omterm{def:b3:topology:compact}{kompak} $K \subseteq \C\setminus\Z$: $P_N \to P$ dan
$P_N' \to P'$ secara seragam
(\cref{thm:b3:holomorphic:weierstrassconv}), dan $\min_K
\abs P > 0$ (sebab $P$ lenyap hanya pada $\Z$), sehingga akhirnya
$\abs{P_N} \geq \frac12\min_K\abs P$ dan $P_N'/P_N \to
P'/P$ secara seragam pada $K$. Sedangkan anggota tengahnya konvergen ke
$\frac1z + \sum_{n\geq1}\frac{2z}{z^2-n^2} = \pi\cot(\pi z)$
menurut pertanyaan 9: sehingga $P'/P = \pi\cot(\pi z)$ pada
$\C\setminus\Z$.

\textbf{13.} Fungsi $\sin(\pi z)$ dan $P$ bersifat utuh dengan himpunan
akar yang sama $\Z$, dan semua akarnya sederhana (pertanyaan 1 dan 11). Di dekat
$m \in \Z$ tulislah $\sin(\pi z) = (z - m)\,\sigma(z)$ dan
$P(z) = (z - m)\,\psi(z)$ dengan $\sigma, \psi$ yang \omterm{def:b3:holomorphic:holo}{holomorfik}
dan tak lenyap di $m$ (faktorkanlah deret pangkatnya): maka $Q =
\sin(\pi z)/P = \sigma/\psi$ diperluas secara \omterm{def:b3:holomorphic:holo}{holomorfik} dan
bebas akar melintasi setiap bilangan bulatnya, dan ia bebas akar pada
$\C\setminus\Z$ sebagai hasil bagi fungsi yang bebas akar. Di sana,
\[
\frac{Q'}{Q} = \frac{(\sin\pi z)'}{\sin\pi z} -
\frac{P'}{P} = \pi\cot(\pi z) - \pi\cot(\pi z) = 0 ,
\]
sehingga fungsi utuh $Q'$ lenyap pada $\C\setminus\Z$,
jadi di mana-mana lewat \omterm{def:b3:topology:continuity}{kekontinuannya}: sehingga $Q$ konstan. Lalu saat $z \to
0$: $\sin(\pi z)/z \to \pi$ dan $P(z)/z \to 1$, sehingga $Q =
\pi$:
\[
\sin(\pi z) = \pi z\prod_{n\geq1}
\Bigl(1 - \frac{z^2}{n^2}\Bigr) .
\]

\textbf{14.} Di $z = \frac12$: $1 = \sin\frac\pi2 =
\frac\pi2\prod_{n\geq1}\bigl(1 - \frac1{4n^2}\bigr) =
\frac\pi2\prod\frac{4n^2-1}{4n^2}$, sehingga
\[
\frac\pi2 = \prod_{n\geq1}\frac{4n^2}{4n^2 - 1}
= \lim_{N\to\infty}\prod_{n=1}^N
\frac{(2n)(2n)}{(2n-1)(2n+1)}
= \lim_{N\to\infty}
\frac{2\cdot2\cdot4\cdot4\cdots(2N)(2N)}
{1\cdot3\cdot3\cdot5\cdots(2N-1)(2N+1)} :
\]
yakni hasil kali Wallis, sebuah akibat satu baris dari pemfaktoran
Euler.

\textbf{15.} Untuk $\abs z \leq r < 1$ setiap faktornya terletak di
$D(1, r^2) \subseteq D(1, 1)$, sehingga $P(z)/z = \exp\bigl(
\sum_{n\geq1}\log(1 - z^2/n^2)\bigr)$: sebab setiap hasil kali parsialnya
adalah eksponensial sebuah jumlah parsial, dan kedua ruasnya beralih ke
limitnya lewat \omterm{def:b3:topology:continuity}{kekontinuan} $\exp$. Sedangkan deret rangkapnya
\[
\sum_{n\geq1}\log\Bigl(1 - \frac{z^2}{n^2}\Bigr)
= -\sum_{n\geq1}\sum_{k\geq1}\frac{z^{2k}}{k\,n^{2k}}
= -\sum_{k\geq1}\frac{\zeta(2k)}k\,z^{2k}
= -\zeta(2)\,z^2 + O(z^4)
\]
tersusun ulang lewat Fubini untuk deret: sebab $\sum_{n,k}
\frac{r^{2k}}{kn^{2k}} \leq \sum_k\zeta(2k)r^{2k} \leq
\zeta(2)\frac{r^2}{1-r^2} < \infty$. Karena itu
\[
P(z) = z\,\exp\bigl(-\zeta(2)z^2 + O(z^4)\bigr)
= z - \zeta(2)\,z^3 + O(z^5) ,
\]
dan pertanyaan 13 membandingkannya dengan $\sin(\pi z) = \pi z -
\frac{\pi^3}6z^3 + O(z^5)$: sehingga $\pi\zeta(2) = \frac{\pi^3}6$,
yakni $\zeta(2) = \frac{\pi^2}6$. Jadi wajah penjumlahannya
(pertanyaan 10) dan wajah perkaliannya menghitung bilangan yang
sama.

\textbf{16.} Pada sebuah \omterm{def:b3:topology:compact}{kompak} $K \subseteq \C\setminus\Z$
jumlah parsialnya $S_N = \frac1z + \sum_{n\leq N}\bigl(
\frac1{z-n} + \frac1{z+n}\bigr)$ (pertanyaan 9, dengan sukunya
dikelompokkan ulang sebagai $\frac{2z}{z^2-n^2} = \frac1{z-n} +
\frac1{z+n}$) konvergen seragam ke $\pi\cot(\pi z)$, sehingga
\cref{thm:b3:holomorphic:weierstrassconv} memberikan $S_N' \to
(\pi\cot\pi z)' = -\pi^2/\sin^2(\pi z)$ secara seragam pada $K$.
Lalu karena $S_N' = -\sum_{\abs n\leq N}(z - n)^{-2}$:
\[
\frac{\pi^2}{\sin^2(\pi z)} =
\sum_{n\in\Z}\frac1{(z - n)^2} ,
\]
dengan kekonvergenannya normal pada \omterm{def:b3:topology:compact}{kompak} $\C\setminus\Z$
(sebab sukunya $O(n^{-2})$).

\textbf{17.} Di $z = \frac12$ ruas kirinya adalah $\pi^2$; sedangkan pada
ruas kanannya, $(\frac12 - n)^2 = \frac{(2n-1)^2}4$ dengan $2n -
1$ melintasi setiap bilangan bulat gasal tepat sekali saat $n$ melintasi
$\Z$:
\[
\pi^2 = \sum_{n\in\Z}\frac{4}{(2n-1)^2}
= 8\sum_{m\geq0}\frac1{(2m+1)^2},
\qquad
\sum_{m\geq0}\frac1{(2m+1)^2} = \frac{\pi^2}8 .
\]
Lalu pemisahan paritasnya: $\zeta(2) = \frac{\pi^2}8 +
\sum_{n\geq1}\frac1{(2n)^2} = \frac{\pi^2}8 +
\frac{\zeta(2)}4$, sehingga $\frac34\zeta(2) = \frac{\pi^2}8$ dan
$\zeta(2) = \frac{\pi^2}6$ sekali lagi.

\textbf{18.} Dengan $c = \cot\theta$ dan $\cot(2\theta) =
\frac{c^2-1}{2c}$: $\cot\theta - 2\cot(2\theta) = c -
\frac{c^2-1}c = \frac1c = \tan\theta$. Karena itu $\pi\tan(\pi z)
= \pi\cot(\pi z) - 2\pi\cot(2\pi z)$, dan pertanyaan 9 di $z$
dan di $2z$ memberikan
\[
\pi\cot(\pi z) = \frac1z +
\sum_{n\geq1}\frac{2z}{z^2 - n^2},
\qquad
2\pi\cot(2\pi z) = \frac1z +
\sum_{n\geq1}\frac{8z}{4z^2 - n^2} .
\]
Kedua deretnya konvergen mutlak di setiap $z$ yang tetap di luar
kutubnya, sehingga selisihnya boleh dikelompokkan ulang sesuka hati: pada
deret keduanya suku genap $n = 2m$ memberikan $\frac{8z}{4z^2 -
4m^2} = \frac{2z}{z^2 - m^2}$ dan meniadakan deret pertamanya
seluruhnya, menyisakan
\[
\pi\tan(\pi z) = -\sum_{m\geq0}\frac{8z}{4z^2 - (2m+1)^2}
= \sum_{m\geq0}\frac{8z}{(2m+1)^2 - 4z^2} ,
\]
secara normal pada \omterm{def:b3:topology:compact}{kompak} yang menghindari $\frac12 + \Z$ (sebab sukunya
$O(m^{-2})$).

\textbf{19.} Untuk $\abs z \leq r < \frac12$:
\[
\frac{8z}{(2m+1)^2 - 4z^2} = \frac{8z}{(2m+1)^2}
\sum_{k\geq0}\Bigl(\frac{4z^2}{(2m+1)^2}\Bigr)^{k},
\]
dengan $\sum_{m,k}8r\,(4r^2)^k(2m+1)^{-2k-2} < \infty$ sebab
$4r^2 < 1$: sehingga Fubini menyusun ulang jumlah rangkapnya menjadi
\[
\pi\tan(\pi z) =
\sum_{k\geq0}8\cdot4^k\,\lambda(2k+2)\,z^{2k+1} .
\]
Lalu terhadap $\pi\tan(\pi z) = \pi^2z + \frac{\pi^4}3z^3 +
O(z^5)$: koefisien $z$-nya memberikan $8\lambda(2) = \pi^2$
--- yakni pertanyaan 17 lagi --- sedangkan koefisien $z^3$-nya memberikan
$32\lambda(4) = \frac{\pi^4}3$, yakni $\lambda(4) =
\frac{\pi^4}{96}$. Lalu dengan membuang penyebut genapnya,
$\lambda(4) = \zeta(4) - 2^{-4}\zeta(4) =
\frac{15}{16}\zeta(4)$: sehingga $\zeta(4) = \frac{16}{15}\cdot
\frac{\pi^4}{96} = \frac{\pi^4}{90}$, yang cocok dengan pertanyaan 7.

\textbf{20.} Berlaku $\cot\frac\theta2 - \cot\theta =
\frac{\cos\frac\theta2\,\sin\theta -
\cos\theta\,\sin\frac\theta2}{\sin\frac\theta2\,\sin\theta}
= \frac{\sin\frac\theta2}{\sin\frac\theta2\,\sin\theta} =
\frac1{\sin\theta}$, sebab pembilangnya adalah $\sin(\theta -
\frac\theta2)$. Lalu dengan $\theta = \pi z$, pertanyaan 9 di
$\frac z2$ berbunyi $\pi\cot\frac{\pi z}2 = \frac2z +
\sum_{n\geq1}\frac{4z}{z^2 - 4n^2}$, sehingga
\[
\frac{\pi}{\sin(\pi z)}
= \pi\cot\frac{\pi z}2 - \pi\cot(\pi z)
= \frac1z + \sum_{n\geq1}\frac{4z}{z^2 - 4n^2}
- \sum_{n\geq1}\frac{2z}{z^2 - n^2} .
\]
Lalu kekonvergenan mutlaknya mengizinkan pengelompokan ulang menurut paritas: sebab $n
= 2m$ yang genap pada deret yang dikurangkan menyumbang
$\frac{2z}{z^2-4m^2}$, menyisakan $+\frac{2z}{z^2-4m^2}$ dari
jumlah pertamanya, sedangkan $n$ yang gasal bertahan bertanda $-$:
\[
\frac{\pi}{\sin(\pi z)} = \frac1z +
\sum_{n\geq1}(-1)^n\,\frac{2z}{z^2 - n^2} .
\]
Untuk pemeriksaan tandanya: $\operatorname{Res}\bigl(\pi/\sin(\pi z),
n\bigr) = \pi/(\pi\cos\pi n) = (-1)^n$
(\cref{met:b3:residues:compute}), dan $(-1)^n\frac{2z}{z^2 -
n^2} = \frac{(-1)^n}{z-n} + \frac{(-1)^n}{z+n}$ membawa
tepat \omterm{def:b3:residues:singularities}{residu} itu di $\pm n$.

\textbf{21.} Untuk $\abs z \leq r < 1$, dengan menguraikan setiap sukunya seperti
pada pertanyaan 19 (yakni $(-1)^n\frac{2z}{z^2-n^2} =
(-1)^{n-1}\frac{2z}{n^2}\sum_k\frac{z^{2k}}{n^{2k}}$) lalu
menerapkan Fubini:
\[
\frac{\pi}{\sin(\pi z)} = \frac1z +
2\sum_{k\geq0}\eta(2k+2)\,z^{2k+1},
\qquad
\eta(s) = \sum_{n\geq1}\frac{(-1)^{n-1}}{n^s} .
\]
Sedangkan pada sisi Taylornya: $\sin(\pi z) = \pi z(1 - \frac{(\pi z)^2}6 +
O(z^4))$ memberikan $\frac\pi{\sin\pi z} = \frac1z +
\frac{\pi^2}6z + O(z^3)$: sehingga $2\eta(2) = \frac{\pi^2}6$, jadi
$\eta(2) = \frac{\pi^2}{12}$. Untuk kekonsistenannya: $\eta(2) =
\zeta(2) - 2\sum_n(2n)^{-2} = (1 - 2^{1-2})\zeta(2) =
\frac{\zeta(2)}2 = \frac{\pi^2}{12}$.

\textbf{22.} Di $z = \frac12$, pertanyaan 20 memberikan
\[
\pi = 2 + \sum_{n\geq1}(-1)^n\frac{1}{\frac14 - n^2}
= 2 + 4\sum_{n\geq1}\frac{(-1)^{n-1}}{(2n-1)(2n+1)} .
\]
Lalu dengan $\frac1{(2n-1)(2n+1)} = \frac12\bigl(\frac1{2n-1} -
\frac1{2n+1}\bigr)$, dan kedua deret berganti tandanya konvergen
(menurut kriteria Leibniz), jumlahnya terbelah menjadi $\frac12\bigl[L -
(1 - L)\bigr] = L - \frac12$, dengan $L = 1 - \frac13 +
\frac15 - \cdots$ dan $\frac13 - \frac15 + \frac17 - \cdots
= 1 - L$. Karena itu $\pi = 2 + 4(L - \frac12) = 4L$:
\[
\frac\pi4 = 1 - \frac13 + \frac15 - \frac17 + \cdots
\]
Moralnya: bahwa setiap kernelnya bersifat \omterm{def:b3:residues:singularities}{meromorfik} dengan kutub pada sebuah
deret aritmetika dan \omterm{def:b3:residues:singularities}{residu} yang ditetapkan. Sedangkan
kotangennya menaruh \omterm{def:b3:residues:singularities}{residu} $1$ di setiap bilangan bulat lalu menjumlahkan
$f(n)$; turunannya mengkuadratkan kutubnya lalu menghargai
$\lambda(2)$; tangennya menggeser kutubnya ke $\frac12 + \Z$
lalu menghargai penyebut gasalnya; sedangkan $\pi/\sin$ menyimpan
kutub bulatnya tetapi menyelang-nyelingkan \omterm{def:b3:residues:singularities}{residunya} $(-1)^n$, sehingga
menghasilkan deret berganti tandanya. Adapun kernel orde pertamanya semuanya
fungsi gasal: sehingga memasangkan $n$ dengan $-n$ menggandakan koefisien
berpangkat genapnya dan memusnahkan yang gasal, jadi $\zeta(2k)$
tercurah secara mekanis sedangkan $\zeta(3)$ tak pernah muncul. Jadi
kebutaannya adalah paritas, bukan kekurangan teknik.

\textbf{23.} Dengan $w = 2\iu\pi z$:
\[
\iu\pi z + \frac{2\iu\pi z}{\eu^{2\iu\pi z} - 1}
= \iu\pi z\,\frac{\eu^{2\iu\pi z} + 1}{\eu^{2\iu\pi z} - 1}
= \iu\pi z\,\frac{\eu^{\iu\pi z} +
\eu^{-\iu\pi z}}{\eu^{\iu\pi z} - \eu^{-\iu\pi z}}
= \pi z\,\frac{\cos\pi z}{\sin\pi z} = \pi z\cot(\pi z) .
\]
Karena itu, dengan memakai fungsi pembangkitnya pada $w = 2\iu\pi z$ (dan
$B_1 = -\frac12$ yang meniadakan suku $\iu\pi z$-nya, sedangkan $B$ gasalnya
lenyap selebihnya):
\[
\pi z\cot(\pi z) = \sum_{k\geq0}\frac{B_{2k}}{(2k)!}
(2\iu\pi z)^{2k}
= 1 + \sum_{k\geq1}(-1)^k\frac{(2\pi)^{2k}B_{2k}}{(2k)!}
z^{2k} .
\]
Sedangkan di sisi lain uraian pecahan parsialnya (Bagian IV)
memberikan $\pi z\cot(\pi z) = 1 - 2\sum_{k\geq1}\zeta(2k)z^{2k}$
(uraikanlah setiap $\frac{2z^2}{z^2 - n^2} =
-2\sum_k\frac{z^{2k}}{n^{2k}}$ lalu jumlahkan atas $n$, dengan kekonvergenan
normalnya membenarkan penukarannya untuk $\abs z < 1$).
Lalu dengan membandingkan koefisiennya: $-2\zeta(2k) =
(-1)^k\frac{(2\pi)^{2k}B_{2k}}{(2k)!}$, yakni rumus yang dinyatakan itu.

\textbf{24.} Untuk $k = 1$: $\frac{(2\pi)^2}{2\cdot2}\cdot\frac16
= \frac{\pi^2}6$. Untuk $k = 2$: $-\frac{(2\pi)^4}{2\cdot24}
\cdot\bigl(-\frac1{30}\bigr) = \frac{16\pi^4}{48\cdot30} =
\frac{\pi^4}{90}$. Dan untuk $k = 3$: $\frac{(2\pi)^6}{2\cdot720}
\cdot\frac1{42} = \frac{64\pi^6}{1440\cdot42} =
\frac{\pi^6}{945}$.

\textbf{25.} Tulis $C(z) = \pi z\cot(\pi z) = 1 -
2\sum_{k\geq1}\zeta(2k)z^{2k}$ (pertanyaan 23). Lalu lewat pendiferensialan
langsung $C = \pi z\cot(\pi z)$, dengan memakai $(\cot u)'
= -1 - \cot^2u$:
\[
zC'(z) = \pi z\cot(\pi z) - \pi^2z^2\bigl(1 +
\cot^2(\pi z)\bigr) = C - \pi^2z^2 - C^2 .
\]
Kini uraikan kedua ruasnya dalam pangkat $z^2$. Ruas kirinya:
$\sum_k(-4k)\,\zeta(2k)\,z^{2k}$. Sedangkan ruas kanannya: dengan mengkuadratkan
deretnya,
\[
C - C^2 = 2\sum_{k\geq1}\zeta(2k)z^{2k}
- 4\sum_{k\geq2}\Bigl(\sum_{j=1}^{k-1}\zeta(2j)
\zeta(2k-2j)\Bigr)z^{2k},
\]
dan suku $-\pi^2z^2 = -6\zeta(2)z^2$ hanya menyesuaikan $k =
1$. Lalu dengan membandingkan koefisien $z^{2k}$ untuk $k \geq 2$:
\[
-4k\,\zeta(2k) = 2\,\zeta(2k) -
4\sum_{j=1}^{k-1}\zeta(2j)\,\zeta(2k-2j),
\]
yakni $\bigl(k + \frac12\bigr)\zeta(2k) =
\sum_{j=1}^{k-1}\zeta(2j)\zeta(2k-2j)$. (Di $k = 1$
identitasnya berbunyi $-4\zeta(2) = 2\zeta(2) - 6\zeta(2)$: yakni sebuah
pemeriksaan kekonsistenan, bukan keterangan baru.) Untuk terapannya: $k =
2$: $\frac52\zeta(4) = \zeta(2)^2 = \frac{\pi^4}{36}$, sehingga
$\zeta(4) = \frac{\pi^4}{90}$; sedangkan $k = 3$: $\frac72\zeta(6) =
2\zeta(2)\zeta(4) = \frac{\pi^6}{270}$, sehingga $\zeta(6) =
\frac{2}{7}\cdot\frac{\pi^6}{270} = \frac{\pi^6}{945}$. Jadi satu
benih transenden ($\zeta(2) = \frac{\pi^2}6$), lalu aljabar murni
menghasilkan setiap nilai zeta genapnya.

\end{solution}
